% VERSAO 2.0: TIPOS, ESCOPOS E STUDIO - 25/11/2026.
% Artigo completo e autocontido: cole TODO este arquivo no main.tex do Overleaf.
% Compilador recomendado no Overleaf: pdfLaTeX. Não requer imagens ou arquivo .bib.
% Identificação da equipe atualizada conforme solicitado.
\documentclass[conference,a4paper]{IEEEtran}
\usepackage[T1]{fontenc}
\usepackage[utf8]{inputenc}
\usepackage[brazil]{babel}
\usepackage{amsmath,amssymb}
\usepackage{array,booktabs,tabularx}
\usepackage{graphicx}
\usepackage{listings}
\usepackage{microtype}
\usepackage{url}
\usepackage[hidelinks]{hyperref}
\urlstyle{same}
\newcolumntype{Y}{>{\raggedright\arraybackslash}X}
\setlength{\emergencystretch}{1.5em}
\lstset{basicstyle=\ttfamily\scriptsize,columns=fullflexible,
  keepspaces=true,showstringspaces=false,breaklines=true,
  frame=single,captionpos=b,aboveskip=6pt,belowskip=6pt,
  xleftmargin=3pt,xrightmargin=3pt}
\addto\captionsbrazil{%
  \renewcommand{\abstractname}{Resumo}%
  \renewcommand{\IEEEkeywordsname}{Palavras-chave}%
  \renewcommand{\refname}{Referências Bibliográficas}%
  \renewcommand{\tablename}{Tabela}%
  \renewcommand{\lstlistingname}{Listagem}%
}
\title{Uma Linguagem Específica de Domínio para\\
Gerenciamento de Estacionamentos Inteligentes}
% Cada autor possui seu próprio bloco de afiliação.
\author{%
  \setlength{\tabcolsep}{22pt}%
  \begin{tabular}{@{}cc@{}}
    Ryan Augusto Dias & Bruno Barroso dos Santos \\
    {\small Curso de Engenharia de Computação} & {\small Curso de Engenharia de Computação} \\
    {\small Fundação Hermínio Ometto} & {\small Fundação Hermínio Ometto} \\
    {\small Campus Araras} & {\small Campus Araras} \\
    {\small Araras, SP, Brasil} & {\small Araras, SP, Brasil}
  \end{tabular}%
}
\begin{document}
\maketitle

\begin{abstract}
Operações de estacionamento envolvem vagas, veículos e transições de ocupação que podem ser representadas por uma linguagem específica de domínio. Nesse recorte, reconhecer comandos textuais não garante que suas sequências respeitem as condições de presença e disponibilidade dos recursos. Este trabalho apresenta uma DSL interpretada para declarar vagas e executar entrada, estacionamento, liberação e saída com validação de estado, aliases tipados, tabela de símbolos, escopos léxicos e interface local. O método compreendeu especificação em EBNF, implementação de lexer e parser manuais em Python, construção de árvore sintática abstrata e avaliação por testes, modelo independente e cargas sintéticas. Foram aprovados 81 métodos de teste e, em cargas de 10, 100 e 1.000 vagas, com cinco aquecimentos e 30 repetições por tamanho, as medianas totais foram 0,390, 3,926 e 41,079 ms, enquanto três pares equivalentes mantiveram a mesma quantidade de linhas e apresentaram razões de caracteres DSL/Python entre 0,676 e 0,787. Conclui-se que o protótipo representou as operações selecionadas e preservou os invariantes nos casos avaliados, com resultados restritos ao ambiente experimental, sem comprovação de usabilidade ou desempenho em estacionamentos reais.
\end{abstract}

\begin{IEEEkeywords}
linguagem específica de domínio; compiladores; interpretação; análise semântica; estacionamento inteligente.
\end{IEEEkeywords}

\section{Introdução}

Estacionamentos inteligentes constituem um domínio que envolve informação sobre disponibilidade, ocupação e utilização de vagas. Lin, Rivano e Le Mouël \cite{lin2017} revisam soluções desse ecossistema, enquanto Geng e Cassandras \cite{geng2013} investigam alocação de recursos e reservas. Dentro desse domínio, o presente trabalho concentra-se na representação textual de operações elementares sobre vagas e veículos. O termo inteligente identifica o contexto da aplicação; o protótipo não utiliza inteligência artificial, sensores ou reconhecimento automático de placas.


A questão delimitada é como expressar uma sequência de operações e verificar sua consistência por meio de uma linguagem própria. Uma instrução pode estar escrita na ordem correta e, ainda assim, solicitar o uso de uma vaga inexistente ou ocupada. Reconhecer comandos não resolve todo o problema: é necessário definir seu significado e as condições em que podem alterar o estado. A justificativa do projeto está em reunir gramática formal, processamento de linguagem e regras de domínio em uma aplicação concreta, sem pressupor a inexistência de sistemas comerciais ou linguagens semelhantes.


O objetivo geral foi desenvolver uma Linguagem Específica de Domínio, ou DSL, interpretada, para representar e executar operações de gerenciamento de estacionamento. Os objetivos específicos foram definir a gramática, implementar reconhecimento de tokens e análise sintática, construir uma árvore sintática abstrata, verificar declaração e ocupação, produzir diagnósticos localizáveis e avaliar programas válidos e inválidos. A questão de pesquisa é: em que medida a DSL representa o conjunto de operações selecionado e preserva os invariantes de vagas e veículos durante a execução?


A implementação adota Python e interpretação por caminhamento de árvore. A linguagem é externa, com arquivos próprios, palavras-chave em português e execução sequencial. A decisão segue um recorte de construção de DSL que exige escolhas explícitas nas etapas de análise, projeto e implementação \cite{mernik2005}. A avaliação compara estados e diagnósticos com expectativas derivadas da especificação e mede o processamento de programas sintéticos de tamanhos controlados.


Os resultados registrados aprovaram os 81 métodos de teste e apresentaram mediana total de 41,079 ms para a carga de 1.000 vagas e 5.000 comandos. A comparação de concisão identificou menor número de caracteres em três cenários da DSL, com a mesma quantidade de linhas do script Python. Esses achados sustentam a contribuição de uma especificação acompanhada de implementação e avaliação reproduzível para o estudo de linguagens; a validade é restrita aos cenários e ao ambiente utilizados, sem estudo com usuários ou avaliação em estacionamento real.


A Seção II apresenta a fundamentação e os trabalhos relacionados. A Seção III descreve a linguagem, a arquitetura e o método experimental. A Seção IV apresenta os resultados e sua análise, e a Seção V reúne a conclusão e as possibilidades de continuidade.


\section{Referencial Teórico e Trabalhos Relacionados}

Uma DSL concentra suas construções em um domínio de aplicação. Van Deursen, Klint e Visser \cite{deursen2000} organizam a literatura sobre essas linguagens e suas estratégias de implementação. Nesta proposta, a especialização aparece nos comandos vaga, entrada, estacionar, liberar e saida. O vocabulário aproxima a notação das operações selecionadas, mas sua presença isolada não demonstra maior usabilidade.


Mernik, Heering e Sloane \cite{mernik2005} discutem padrões nas fases de decisão, análise, projeto e implementação de DSLs. Essa organização fundamenta a separação entre necessidades do domínio e escolhas sobre a linguagem. A construção de ferramentas próprias também introduz custo de desenvolvimento e manutenção. Por esse motivo, o escopo reúne cinco operações físicas, declarações de aliases tipados e blocos léxicos, sem funções, laços ou expressões de propósito geral.


No domínio de aplicação, Lin et al. \cite{lin2017} apresentam um levantamento de soluções de estacionamento inteligente, considerando funcionalidades e tecnologias associadas. O estudo contextualiza o problema, mas não fornece a especificação da DSL aqui implementada. O protótipo utiliza apenas a camada lógica: eventos são fornecidos por programas textuais, sem aquisição física de informações.


Geng e Cassandras \cite{geng2013} propõem alocação e reserva de vagas com base em uma função de custo do motorista e otimização de recursos. No presente projeto, a vaga é indicada explicitamente pelo programa, sem alocação ótima, reserva ou precificação. O interesse comum está na representação consistente de recursos; as finalidades e os critérios de avaliação são distintos.


Como abordagem recente de engenharia de linguagens, Stojanovi\'{c}, Kova\v{c}evi\'{c} e Kosar \cite{domainbuilder2026} apresentam DomainBuilder, um framework em C++ com interpretador extensível e sintaxe baseada em S-expressões. Seu caso de uso envolve pipelines de dados e combina operadores de domínio com construções gerais. A Tabela~\ref{tab:t1} compara características desse framework com o presente protótipo: aqui, a sintaxe própria e as construções de domínio delimitam o escopo acadêmico, sem estabelecer vantagem de desempenho entre ferramentas com finalidades e cargas diferentes.


\begin{table*}[!t]
\centering
\footnotesize
\setlength{\tabcolsep}{4pt}
\renewcommand{\arraystretch}{1.15}
\caption{Comparação de características das abordagens de implementação}
\label{tab:t1}
\begin{tabularx}{\textwidth}{@{}p{0.17\textwidth}YYY@{}}
\toprule
\textbf{Abordagem} & \textbf{Sintaxe e execução} & \textbf{Construções gerais} & \textbf{Validação do estacionamento} \\
\midrule
DomainBuilder \cite{domainbuilder2026} & S-expressões; interpretador extensível em C++. & Sim; combinadas com operadores do domínio. & Não é o caso de uso apresentado. \\
\textbf{Protótipo deste trabalho} & \textbf{Comandos em português; AST e interpretador em Python.} & \textbf{Cinco operações, aliases tipados e blocos.} & \textbf{Tipos, escopos, presença e ocupação.} \\
\bottomrule
\end{tabularx}
\end{table*}

ParkLang, de Ozkaya, Turunc e Cobanoglu~\cite{parklang2026}, é uma linguagem do mesmo domínio, voltada à modelagem arquitetural de sistemas de estacionamento. O trabalho apresenta uma notação gráfica baseada em níveis do modelo C4 e suporte de ferramenta. Diferentemente desse propósito, EstacionaDSL interpreta comandos textuais e verifica presença e ocupação. Portanto, o diferencial aqui é o nível operacional executável e a exposição didática do processamento de linguagem, não uma alegação de inexistência de DSLs de estacionamento ou superioridade entre ferramentas. Não se realizou execução comparativa de ParkLang.

As referências clássicas são empregadas para conceitos estabelecidos e contextualização do domínio, enquanto o trabalho de 2026 situa uma abordagem contemporânea de implementação. A seleção não constitui revisão sistemática. A comparação justifica manter uma gramática própria e separar estrutura de estado: essa decisão torna explícitas as etapas estudadas na disciplina e delimita as regras que serão avaliadas, orientando os materiais e métodos a seguir.


\subsection{Fundamentos de reconhecimento e interpretação}
As classes de identificadores e literais são descritas por expressões regulares. Para um identificador, um autômato conceitual exige uma letra inicial e aceita letras, dígitos ou sublinhado nas transições seguintes. O lexer implementa varredura e correspondência desses padrões, registrando posições. O parser descendente recursivo reconhece comandos por seus tokens iniciais distintos e utiliza a pilha de chamadas nos blocos aninhados. A gramática abstrata admite recursão; a implementação restringe a profundidade a 32 blocos para limitar recursos.

A AST representa construções relevantes, independentemente de espaços e quebras de linha. A interpretação percorre essa estrutura e consulta ambientes de símbolos. Reconhecimento sintático, compatibilidade de tipos e regras de ocupação são problemas distintos: uma entrada pode ter sintaxe correta e referenciar uma vaga onde se exige veículo. A escolha de uma DSL explicita essas restrições e reduz parte da sintaxe incidental de scripts gerais, ao custo de desenvolver e manter gramática, diagnósticos e ferramentas. A evidência de concisão não substitui uma avaliação humana dessa escolha.

\section{Materiais e Métodos}

\subsection{Delineamento, ferramentas e escopo}

A pesquisa é aplicada, com desenvolvimento de protótipo e avaliação experimental. A implementação utiliza Python 3, com biblioteca padrão para expressões regulares, estruturas de dados, testes, serialização e medição de tempo \cite{python}. Foram adotados lexer manual e parser descendente recursivo, mantendo explícitas as etapas do processamento e dispensando geradores de parser.


O programa recebe um arquivo .dsl em UTF-8, produz diagnósticos e apresenta o estado final no terminal. Também permite inspecionar tokens e AST e exportar resultados em JSON. As estruturas permanecem em memória durante uma execução; cada programa inicia com estado independente. O código emprega recursos de Python 3.10 ou superior, e a execução avaliada utilizou CPython 3.15.2.


São contempladas declaração de vagas, entrada, estacionamento, liberação e saída. As categorias comum, pcd e idoso são atributos das vagas. Como não há credenciais ou perfis de condutores na sintaxe, não ocorre verificação de elegibilidade por categoria. Cobrança, reservas, persistência, sensores, integração com serviços externos e operações simultâneas ficaram fora desta versão.


\subsection{Sintaxe concreta e regras léxicas}

Um programa define um estacionamento e contém pelo menos um comando entre chaves. A Listagem~\ref{lst:l1} apresenta um ciclo completo. Liberar desfaz a associação com a vaga, mas mantém o veículo presente; somente saida encerra a presença no estacionamento.


\begin{lstlisting}[caption={Programa válido utilizado na avaliação.},label={lst:l1}]
estacionamento Central {
vaga A1 comum
vaga A2 pcd
vaga A3 idoso
entrada "ABC1234"
estacionar "ABC1234" em A1
liberar A1
saida "ABC1234"
}
\end{lstlisting}

A Listagem~\ref{lst:l2} apresenta o subconjunto original da EBNF, estendido adiante com aliases tipados e blocos. Chaves fora de aspas representam repetição; chaves entre aspas são delimitadores literais. ID designa identificador, PLACA um literal de veículo e EOF o término da entrada.


\begin{lstlisting}[caption={Subconjunto original da gramática em EBNF.},label={lst:l2}]
programa = "estacionamento", ID, "{",
  comando, { comando }, "}", EOF ;
comando = declaracao | entrada | estacionar
  | liberar | saida ;
declaracao = "vaga", ID, tipo ;
tipo = "comum" | "pcd" | "idoso" ;
entrada = "entrada", PLACA ;
estacionar = "estacionar", PLACA,
  "em", ID ;
liberar = "liberar", ID ;
saida = "saida", PLACA ;
\end{lstlisting}

Os identificadores seguem [A-Za-z][A-Za-z0-9\_]*. Após reconhecer o lexema completo, o lexer consulta o conjunto de palavras reservadas; assim, vagaExtra permanece um único ID, e palavras-chave não podem ser nomes. A linguagem distingue maiúsculas e minúsculas e aceita palavras-chave apenas na forma indicada. Espaços, tabulações e quebras de linha são ignorados entre tokens; não há ponto e vírgula nem comentários.


PLACA aceita exatamente sete caracteres de [A-Z0-9] entre aspas duplas. Trata-se de identificador sintético de veículo, não de validação normativa de placas reais. Aspas sem fechamento na mesma linha, caracteres proibidos e literais fora do padrão produzem erro léxico. Cada token armazena tipo, valor, linha e coluna, com posições iniciadas em 1; tabulação conta como um caractere.


\subsection{Arquitetura e processamento da linguagem}

O fluxo é código-fonte, tokens, AST completa, validações e execução. A implementação separa lexer.py, parser.py, ast\_nodes.py, semantic.py, interpreter.py, diagnostics.py, symbols.py, service.py e main.py no pacote parkingdsl. O parser constrói nós DeclararVaga, Entrada, Estacionar, Liberar, Saida, Binding e Block, preservando a posição inicial de cada comando.


A árvore completa é construída antes de executar qualquer comando. Consequentemente, falhas léxicas e sintáticas impedem toda a execução. A sequência estacionar em A1 "ABC1234" é rejeitada pelo parser. A categoria premium também é rejeitada sintaticamente, pois as categorias válidas são enumeradas na gramática; não é tratada como erro semântico de ocupação.


As alternativas de comando têm conjuntos FIRST disjuntos: cada uma começa com palavra-chave diferente. Um token de antecipação basta para selecionar a produção, sem retrocesso. A repetição de comandos utiliza laço, evitando profundidade de pilha proporcional ao programa. Na versão 2.0, blocos locais podem se aninhar; a profundidade da pilha acompanha esse aninhamento, limitado a 32, e não o número de comandos irmãos.


As validações semânticas dependentes do estado são aplicadas na ordem dos comandos, imediatamente antes de cada alteração. Declarações duplicadas e referências a vagas ainda não declaradas seguem essa mesma ordem. Uma declaração posterior não torna válido um uso anterior. Ao encontrar falha semântica, o interpretador interrompe o programa: o comando rejeitado não altera o estado, os efeitos anteriores permanecem e não há reversão integral do programa.


\subsection{Modelo de estado e semântica operacional}

O estado é S = (V, P, O): V mapeia vagas para categorias, P é o conjunto de veículos presentes e O é uma função parcial de vagas ocupadas para veículos. O deve ser injetiva. São preservados os invariantes:


\begin{equation}\operatorname{dom}(O)\subseteq\operatorname{dom}(V),\quad\operatorname{im}(O)\subseteq P.\label{eq:inv}\end{equation}

O código usa dicionários e conjuntos, com um índice inverso I de veículo para vaga. Além de (1) e da injetividade, verifica-se I = O\textsuperscript{-1}. Consultar esse índice permite detectar um veículo já estacionado sem percorrer todas as vagas. As duas direções da associação são alteradas somente após validar todas as pré-condições.


\begin{table*}[!t]
\centering
\footnotesize
\setlength{\tabcolsep}{4pt}
\renewcommand{\arraystretch}{1.15}
\caption{Semântica dos comandos: x é vaga, p é veículo}
\label{tab:t2}
\begin{tabularx}{\textwidth}{@{}p{0.18\textwidth}YY@{}}
\toprule
\textbf{Comando} & \textbf{Pré-condição} & \textbf{Efeito} \\
\midrule
vaga x c & x não declarada. Categoria reconhecida pelo parser. & V[x] = c; vaga livre. \\
entrada p & p ausente. & Adiciona p a P. \\
estacionar p em x & x declarada e livre; p presente e sem vaga. & O[x] = p; I[p] = x. \\
liberar x & x ocupada. & Remove O[x] e I[p]; mantém p em P. \\
saida p & p presente e sem vaga. & Remove p de P. \\
\bottomrule
\end{tabularx}
\end{table*}

A saída exige liberação prévia: o interpretador não libera uma vaga silenciosamente. Em estacionar, a ordem fixa de diagnóstico verifica existência da vaga, presença do veículo, disponibilidade da vaga e ausência de outra vaga para o veículo. Em liberar, distingue-se vaga inexistente de vaga declarada já livre, refinando o diagnóstico sem alterar a condição de validade.


Cada erro informa fase, código, linha, coluna e descrição. Erros semânticos apontam o início do comando; erros sintáticos apontam o token inesperado; erros léxicos apontam o caractere ou a aspa inicial. Somente comandos aceitos geram eventos. A saída apresenta as vagas em ordem lexicográfica determinística e distingue sucesso de interrupção.


A opção --verify confere os invariantes após cada comando aceito. A atomicidade descrita refere-se às rejeições semânticas previstas, não a interrupções externas, falta de memória ou falhas de hardware. O protótipo é sequencial e não implementa controle de concorrência.


\subsection{Tipos, tabela de símbolos e escopos}
Além dos comandos físicos, a versão 2.0 admite aliases imutáveis e blocos locais. A extensão da EBNF é apresentada a seguir; entrada, saída e o primeiro operando de estacionar passam a aceitar PLACA ou ID. O destino e a liberação continuam utilizando ID.
\begin{lstlisting}[caption={Extensão da gramática para tipos e escopos.}]
operando = PLACA | ID ;
definicao = "definir", ID, ":", tipo_dado,
            "=", operando ;
tipo_dado = "veiculo" | "vaga" ;
bloco = "bloco", "{", { comando }, "}" ;
\end{lstlisting}
As alternativas de comando incluem definicao e bloco. O programa exige pelo menos um comando; blocos internos podem ser vazios. Categorias comum, pcd e idoso são atributos enumerados de vagas, não substituem os tipos de dados veiculo e vaga.

Cada entrada na tabela registra nome, tipo, valor, caminho de escopo e posição de declaração. O ambiente é uma pilha de dicionários: a resolução procura primeiro no escopo atual e depois nos externos. Se $\Gamma(x)=\tau$, o identificador possui tipo $\tau$; um literal PLACA possui tipo veiculo. Uma definição anotada $x:\tau=e$ exige que o tipo resolvido de $e$ seja igual a $\tau$. Entrada e saída exigem veiculo; estacionar exige veiculo e vaga, respectivamente. Não há conversão implícita.

Nomes não podem ser repetidos no mesmo ambiente nem usados antes da declaração. Um ambiente interno pode sombrear nomes externos. Ao sair do bloco, sua tabela é removida, inclusive em interrupção por erro; os efeitos físicos já aceitos continuam. Vagas físicas só podem ser declaradas globalmente, mantendo sua identidade durante a execução. Os tipos são checados sequencialmente antes de cada comando, e não em uma passagem estática completa anterior à interpretação.

\begin{lstlisting}[caption={Alias global e sombreamento local.}]
estacionamento Campus {
  vaga A1 comum
  definir cliente: veiculo = "ABC1234"
  entrada cliente
  bloco {
    definir cliente: veiculo = "DEF5678"
    entrada cliente
  }
  estacionar cliente em A1
}
\end{lstlisting}
O exemplo associa A1 a ABC1234 e mantém DEF5678 presente sem vaga. T001 indica incompatibilidade de tipo, T002 nome não declarado ou fora do escopo, T003 duplicidade no mesmo ambiente e T004 declaração física local. O diagnóstico de categoria inválida permanece sintático. Destinos físicos não declarados preservam E002 por compatibilidade com a versão original.

\subsection{Interface gráfica local}
O Studio fornece editor, exemplos, abertura e salvamento de arquivos, execução, mapa de vagas e inspeção de tokens, AST, eventos e símbolos. HTML, CSS e JavaScript compõem a interface; um servidor da biblioteca padrão Python escuta somente no endereço de loopback. O código enviado segue para o mesmo lexer, parser e interpretador da CLI. Não existe um segundo interpretador em JavaScript.

Cada chamada utiliza estado novo. O painel exibe o resultado da última execução e avisa quando o texto foi alterado. O serviço limita fontes a 64 KiB e 6.000 tokens; cargas maiores são executadas pela CLI. O editor não executa código Python e o servidor valida origem e caminhos. Não há hospedagem pública, autenticação ou persistência de estado. A interface apoia demonstrações, sem acrescentar cobrança ou funcionalidades externas ao domínio especificado.

\subsection{Avaliação de corretude e diagnóstico}

Os testes foram organizados separadamente do código de produção, com expectativas codificadas a partir da especificação. Programas válidos verificam estado final e operações aceitas. Nos inválidos, a simples rejeição não é suficiente: os testes de diagnóstico comparam fase, código e posição, enquanto os casos semânticos verificam ausência de mutação pelo comando rejeitado. A integração verifica que não há comandos posteriores executados e que erros estruturais não produzem efeitos.


A bateria contém famílias léxica, sintática, semântica, de programas válidos, de integração e de modelo. Inclui literal inválido, aspas abertas, caractere proibido, chave ou palavra em ausente, categoria inválida, duplicidade, ocupação conflitante, ausência de veículo, liberação de vaga livre e saída sem liberação. Os testes também examinam reentrada, realocação e consistência do índice inverso.


Para ampliar as combinações, um método de teste utiliza um modelo independente baseado em lista de vagas e busca linear, sem chamar o validador ou consultar o índice inverso do interpretador. São usadas dez sementes fixas, com 120 tentativas de operação por semente; somente as candidatas válidas são incorporadas à sequência. Após cada operação incorporada, presença e ocupação são comparadas com o modelo, e o interpretador confere seus invariantes.


A métrica principal é a proporção de métodos de teste aprovados, reportada também por família. Um método pode conter subcasos; estes não são contados como testes independentes. Cada execução usa estado independente. A bateria finita fornece evidência de conformidade aos casos avaliados, sem constituir prova formal de corretude.


\subsection{Desempenho, concisão e reprodutibilidade}

O gerador produz programas de 10, 100 e 1.000 vagas com um ciclo completo por vaga. Cada ciclo contém declaração, entrada, estacionamento, liberação e saída, totalizando 5n comandos para n vagas. Os veículos têm identificadores únicos de sete dígitos, compatíveis com o léxico. A ordem da carga permanece idêntica entre as repetições de cada configuração.


Foram realizadas cinco execuções de aquecimento e trinta repetições medidas por tamanho, sempre com estado novo. As medições usam perf\_counter\_ns e separam tokenização, parsing e interpretação. A geração, a leitura, a impressão, o descarte final dos objetos e a conferência global de invariantes ficam fora do intervalo. A interpretação medida inclui criação do estado, pré-condições semânticas, mutações e registro de eventos; o coletor de lixo padrão permanece habilitado.


A conferência global após cada comando foi desabilitada na medição para não acrescentar varreduras sucessivas do estado. As regras semânticas continuam ativas, e os invariantes são conferidos ao final, fora do intervalo, juntamente com a ausência de veículos presentes e vagas ocupadas. São reportadas mediana e percentil 95, calculado por nearest-rank: amostra na posição ceil(0,95N) da série ordenada.


A concisão compara três cenários, incluindo a Listagem~\ref{lst:l1}, com scripts Python equivalentes apoiado em adaptador para o mesmo núcleo semântico. Contam-se linhas não vazias e não comentadas e caracteres, com formatação fixa, separação LF entre linhas e sem LF final. A contagem inclui delimitadores na DSL e importação e criação de objeto em Python. O tamanho da infraestrutura é informado separadamente. A razão de linhas é:


\begin{equation}R_L=\frac{L_{\mathrm{DSL}}}{L_{\mathrm{Python}}}.\label{eq:conc}\end{equation}

Esse indicador não mede legibilidade, aprendizagem ou produtividade humana. O adaptador Python compartilha o núcleo e não é uma implementação independente para comparar desempenho. Estado final e tipos de operação dos dois scripts são verificados antes de registrar a contagem.


O material suplementar está disponível no repositório público EstacionaDSL~\cite{estacionadsl} e reúne código-fonte, EBNF, exemplos, testes, gerador e resultados em JSON e CSV. Os comandos python run\_tests.py, python benchmark.py e python compare\_concision.py reproduzem os experimentos. Os registros incluem versão de Python, sistema, processador, memória disponível e SHA-256 dos módulos de produção e do benchmark; a execução não depende de pacotes externos à biblioteca padrão.


\section{Resultados e Análise}

\subsection{Implementação e resultados de corretude}

O protótipo implementou as cinco formas de comando e os três atributos de categoria previstos. A execução da Listagem~\ref{lst:l1} aceitou sete comandos e terminou com A1, A2 e A3 livres, conjunto P vazio e associações O e I vazias. Cenários adicionais mantiveram múltiplos veículos estacionados, permitiram reentrada após saída e realocação após liberação.


Os 81 métodos de teste foram aprovados, sem falhas ou erros na execução registrada. A Tabela~\ref{tab:t3} apresenta a distribuição por família. O método de modelo reúne as dez sementes e suas sequências; não representa apenas uma operação isolada. A taxa de 100\% é relativa à bateria executada e não deve ser interpretada como cobertura integral de todos os programas possíveis.


\begin{table}[!t]
\centering
\footnotesize
\setlength{\tabcolsep}{4pt}
\renewcommand{\arraystretch}{1.15}
\caption{Resultados dos métodos de teste}
\label{tab:t3}
\begin{tabularx}{\columnwidth}{@{}Yrrr@{}}
\toprule
\textbf{Família} & \textbf{Total} & \textbf{Aprovados} & \textbf{Taxa} \\
\midrule
Léxica & 9 & 9 & 100\% \\
Sintática & 10 & 10 & 100\% \\
Semântica e invariantes & 14 & 14 & 100\% \\
Programas válidos & 7 & 7 & 100\% \\
Integração e CLI & 9 & 9 & 100\% \\
Modelo independente & 1 & 1 & 100\% \\
Tipos, escopos e serviço & 21 & 21 & 100\% \\
HTTP local & 10 & 10 & 100\% \\
Total & 81 & 81 & 100\% \\
\bottomrule
\end{tabularx}
\end{table}

Os casos de invariantes confirmaram a rejeição de veículo em duas vagas, liberação de vaga livre e saída sem liberação. O cenário de índice inverso inconsistente foi detectado pela conferência de invariantes. Nas sequências válidas geradas, o estado do interpretador coincidiu com o modelo independente após cada operação incorporada, oferecendo evidência adicional além dos exemplos fixos.


\subsection{Diagnósticos e preservação de estado}

Na tentativa de estacionar um segundo veículo em vaga ocupada, o exemplo 04 produziu E005 na linha 6, coluna 1. Os quatro comandos anteriores permaneceram aceitos: declaração de A1, entradas de ABC1234 e DEF5678 e estacionamento do primeiro veículo. O estado parcial manteve O[A1] = ABC1234 e ambos os veículos presentes. A liberação escrita após a falha não foi executada.


No exemplo 05, a ausência da palavra em foi identificada como S001 na linha 4, coluna 22. Apesar de o arquivo conter declaração e entrada em linhas anteriores, nenhuma operação foi aplicada, pois o parsing completo falhou. No exemplo 06, o literal com letras minúsculas gerou L003 na linha 3, coluna 9, também sem efeitos no estado. Já o exemplo 07 rejeitou saída de veículo estacionado com E008 na linha 5, coluna 1, preservando a associação existente.


Esses resultados distinguem validade estrutural e validade operacional. A categoria premium foi rejeitada pelo parser, enquanto conflitos de ocupação foram diagnosticados durante a interpretação. A ordem fixa das verificações também foi exercitada quando mais de uma condição poderia falhar, evitando mensagens dependentes da ordem interna de consulta aos dicionários.


\subsection{Desempenho observado}
Os ensaios da versão 2.0 foram realizados em 25 de novembro de 2026, em Windows 11, CPython 3.15.2, Intel Core i7-9700K e oito processadores lógicos, com 31,92 GiB de memória física. Foram mantidos cinco aquecimentos e 30 repetições por tamanho. Os resultados atualizados substituem os da versão inicial.
\begin{table}[!t]
\centering\footnotesize
\caption{Cargas e tempos totais do núcleo na versão 2.0}
\label{tab:t4}
\begin{tabularx}{\columnwidth}{@{}Yrrr@{}}\toprule
Vagas/comandos & Bytes & Mediana (ms) & p95 (ms) \\\midrule
10 / 50 & 892 & 0,390 & 0,416 \\
100 / 500 & 8905 & 3,926 & 3,959 \\
1000 / 5000 & 91708 & 41,079 & 44,962 \\
\bottomrule\end{tabularx}
\end{table}
A medição compreende lexer, parser e interpretação, incluindo tabela de símbolos, validações e eventos. Exclui interface, comunicação HTTP, ordenação e impressão dos resultados. Não houve estudo da latência visual do painel. As estatísticas detalhadas por etapa e todas as amostras estão em benchmark.json e benchmark.csv.

O crescimento nos três pontos caracteriza as cargas escolhidas, sem provar complexidade assintótica ou garantias de tempo real. A tokenização percorre caracteres, o parsing percorre tokens e a resolução de nomes consulta até a profundidade de escopo, limitada a 32. Consultas a tabelas hash têm custo médio constante para chaves curtas, não pior caso garantido. A verificação global dos invariantes após cada comando é desabilitada no benchmark, mas mantida nos testes e no painel.

As cargas medidas usam ciclos físicos e escopo global. Exercitam o custo da tabela de símbolos de vagas, mas não quantificam o impacto de diferentes profundidades de blocos. Tokens, AST, símbolos e eventos permanecem em memória. Não se mediu pico de memória nem desempenho relativo a ParkLang, DomainBuilder ou outra implementação.

\subsection{Concisão e infraestrutura}
A comparação foi ampliada para três pares fixos: ciclo completo, ocupação simultânea e realocação. Verificou-se igualdade de estado final e tipos de operação antes da contagem. A Tabela~\ref{tab:t6} mostra que a quantidade de linhas foi igual em cada par; os caracteres foram menores na DSL sob a formatação escolhida.
\begin{table}[!t]
\centering\footnotesize
\caption{Três cenários equivalentes: DSL/Python}
\label{tab:t6}
\begin{tabularx}{\columnwidth}{@{}Yrrr@{}}\toprule
Cenário & Linhas & Caracteres & Razão \\\midrule
Ciclo & 9/9 & 138/204 & 0,676 \\
Ocupação & 11/11 & 200/254 & 0,787 \\
Realocação & 11/11 & 181/239 & 0,757 \\
\bottomrule\end{tabularx}
\end{table}
A infraestrutura parkingdsl contém 584 linhas físicas, e o adaptador Python acrescenta 27. O adaptador utiliza o mesmo núcleo semântico e não constitui implementação independente; o modelo independente citado nos testes é outro artefato. A comparação não inclui o tamanho da interface e não avalia a concisão dos novos aliases ou blocos. Os três cenários continuam sendo uma amostra pequena e selecionada pelos autores.

As razões descrevem tamanho de scripts e sofrem influência da formatação e das importações. Não demonstram menor dificuldade, tempo de aprendizagem ou produtividade. Não foi realizado estudo com participantes. A justificativa da DSL é expressar operações e restrições de forma explícita; a medição fornece evidência descritiva limitada dessa escolha.

\subsection{Limitações e ameaças à validade}

A validade interna é limitada pela possibilidade de erros nos testes e nos resultados esperados. A separação entre produção e testes, os estados esperados derivados da especificação e o modelo independente reduzem esse risco, mas não o eliminam. A conferência dos invariantes após comandos nos testes reforça a inspeção do estado intermediário; ainda assim, não equivale a uma demonstração formal para todas as entradas.


Na avaliação de desempenho, aquecimento, repetições e preservação da carga controlam parte da variabilidade. Porém, não houve isolamento experimental do sistema operacional, múltiplas máquinas ou múltiplas versões de Python. As repetições foram realizadas em um único ambiente; percentis e medianas descrevem essas amostras, sem representar garantias universais de latência.


A validade externa é restrita por cargas sintéticas de até 1.000 vagas e por estado inteiramente em memória. Não foram avaliados eventos concorrentes, falhas de infraestrutura, persistência ou operação em estacionamento real. As categorias de vaga não incluem credenciais; a DSL não deve ser interpretada como mecanismo de fiscalização de acessibilidade ou prioridade.


A validade de construto da concisão também é limitada: linhas e caracteres não medem compreensão, aprendizagem ou produtividade. A seleção bibliográfica contextualiza as decisões, sem revisar exaustivamente todas as DSLs ou soluções de estacionamento. Essas limitações orientam a interpretação da contribuição como protótipo acadêmico reproduzível.


A revisão medida é identificada em benchmark.json por SHA-256 dos nomes e bytes dos módulos parkingdsl/*.py e benchmark.py. O valor completo, apresentado em duas partes contíguas na Listagem~\ref{lst:l3}, permite conferir a correspondência entre o código e as medições; não identifica, sozinho, todos os arquivos de documentação e teste.


\begin{lstlisting}[caption={SHA-256 da revisão medida; concatenam-se as duas linhas.},label={lst:l3}]
f350938cd9ca4ee38013e137f79ea143
3dbe768c087889b1f83c59cef201f456
\end{lstlisting}

\section{Conclusão}

Este trabalho especificou e implementou uma DSL externa e interpretada para operações elementares de estacionamento. A solução reúne gramática, lexer, parser descendente recursivo, AST, regras semânticas e interpretação sequencial em Python. A separação entre reconhecimento estrutural e validação de estado permitiu produzir diagnósticos localizáveis e definir explicitamente os efeitos de interrupções.


Em resposta à questão de pesquisa, a DSL representou o conjunto de cinco operações selecionado e preservou os invariantes nos cenários avaliados. Os 81 métodos de teste foram aprovados, incluindo estados parciais, ausência de efeitos de comandos inválidos e comparação com modelo independente. Essa evidência sustenta o atendimento aos casos de aceitação implementados, sem constituir prova formal de corretude.


A avaliação caracterizou o custo do processamento neste ambiente e identificou a tokenização como etapa de maior tempo observado. A análise de concisão mostrou economia de caracteres em três cenários, sem ganho de linhas nem comprovação de usabilidade. Assim, os resultados oferecem uma base verificável para demonstrações e exercícios de processamento de linguagens, sem qualificar o protótipo como sistema operacional de estacionamento.


A contribuição final é uma implementação acompanhada de especificação e evidências reproduzíveis para estudo de processamento de linguagens. Como continuidade, podem ser investigados persistência, perfis de condutores, novas operações e tratamento de concorrência, cada qual acompanhado de revisão da semântica e de novos testes. Avaliações com usuários e medições em outras máquinas e cargas são necessárias antes de formular conclusões sobre facilidade de uso, robustez operacional ou aplicação real.


Os principais desafios de implementação foram separar presença de ocupação, preservar o prefixo válido em erros semânticos e restaurar ambientes após blocos. Foram tratados por índices inversos, validação antes de mutação e remoção de escopo em finally. A revisão acrescentou tipos e escopos explícitos e tornou o processamento observável na interface. Como evolução experimental, são necessários participantes reais e cenários externos aos elaborados pelos autores antes de afirmar ganhos de usabilidade.

\begin{thebibliography}{00}

\bibitem{lin2017}
T. Lin, H. Rivano e F. Le Mouël, ``A survey of smart parking solutions,'' \emph{IEEE Transactions on Intelligent Transportation Systems}, vol. 18, n. 12, pp. 3229-3253, 2017. doi: \href{https://doi.org/10.1109/TITS.2017.2685143}{10.1109/TITS.2017.2685143}.

\bibitem{geng2013}
Y. Geng e C. G. Cassandras, ``New `smart parking' system based on resource allocation and reservations,'' \emph{IEEE Transactions on Intelligent Transportation Systems}, vol. 14, n. 3, pp. 1129-1139, 2013. doi: \href{https://doi.org/10.1109/TITS.2013.2252428}{10.1109/TITS.2013.2252428}.

\bibitem{mernik2005}
M. Mernik, J. Heering e A. M. Sloane, ``When and how to develop domain-specific languages,'' \emph{ACM Computing Surveys}, vol. 37, n. 4, pp. 316-344, 2005. doi: \href{https://doi.org/10.1145/1118890.1118892}{10.1145/1118890.1118892}.

\bibitem{deursen2000}
A. van Deursen, P. Klint e J. Visser, ``Domain-specific languages: An annotated bibliography,'' \emph{ACM SIGPLAN Notices}, vol. 35, n. 6, pp. 26-36, 2000. doi: \href{https://doi.org/10.1145/352029.352035}{10.1145/352029.352035}.

\bibitem{domainbuilder2026}
A. Stojanovi\'{c}, \v{Z}. Kova\v{c}evi\'{c} e T. Kosar, ``Design and implementation of a domain-specific language framework using S-expressions,'' \emph{Journal of Computer Languages}, vol.~87, art.~101398, 2026, doi: \href{https://doi.org/10.1016/j.cola.2026.101398}{10.1016/j.cola.2026.101398}.

\bibitem{parklang2026}
M. Ozkaya, A. Turunc e U. Cobanoglu, ``Towards an architecture modeling language for smart car parking management systems,'' em \emph{Proceedings of MODELSWARD}, pp. 331--340, 2026. doi: \href{https://doi.org/10.5220/0014399800004058}{10.5220/0014399800004058}.

\bibitem{python}
Python Software Foundation, ``The Python standard library,'' \emph{Python 3.15.2 documentation}. [Online]. Disponível em: \href{https://docs.python.org/3/library/}{https://docs.python.org/3/library/}. Acesso em: 23 nov. 2026.


\bibitem{estacionadsl}
R. A. Dias e B. B. dos Santos, ``EstacionaDSL: código-fonte, exemplos e resultados do projeto,'' \emph{GitHub}, 2026. [Online]. Disponível em: \url{https://github.com/01ryu/estaciona-dsl}. Acesso em: 25 nov. 2026.

\end{thebibliography}
\end{document}
